\newpage
\section{Periodic Cohomology}
\subsection{Group Cohomology of a Cyclic Group}
Now we will compute the cohomology of a group in a simpler case, when the group is cyclic.
Let $G$ be a cyclic group of order $n$ generated by $\sigma$.

Before beginning the computation, we observe the following:
\begin{remark}
    \noindent
    \begin{enumerate}
        \item $\Z[\Z] = \Z[t, t^{-1}]$.
        \item $\Z[\Z_{>0}] = \Z[t]$.
        \item $\Z[\Z/n\Z] = \Z[t]/(t^n-1)$.
    \end{enumerate}
\end{remark}
So we now know that the group ring $\Z[G]$ is isomorphic to the quotient ring $\Z[t]/(t^n-1)$
via the map $\sigma \mapsto t$.
Let $N$ be the norm element $N = 1 + \sigma + \cdots + \sigma^{n-1} \in \Z[G]$.
In $\Z[t]/(t^n-1)$, we have $N = 1 + t + \cdots + t^{n-1}$.
So we observe that $N(\sigma-1)$ is corresponding to $N(t-1) = (1+t+\cdots+t^{n-1})(t-1) = t^n - 1 = 0$ in $\Z[t]/(t^n-1)$.
It follows that $\im(N) \subseteq \ker(\sigma-1)$
and $\im(\sigma-1) \subseteq \ker(N)$
when we consider $N$ and $\sigma-1$ as multiplication maps on $\Z[G]$.
Moreover, we have
\begin{proposition}
    $\ker(N) = \im(\sigma-1)$ and $\ker(\sigma-1) = \im(N)$.
\end{proposition}
\begin{proof}
    Suppose that $a \in \ker(N)$, i.e., $Na = 0$.
    Let $a(t)$ be the polynomial in $\Z[t]$ corresponding to $a$.
    Then, we have $N(t)a(t) = (t^n-1)b(t)$ for some $b(t) \in \Z[t]$.
    Since $N(t)(t-1) = t^n-1$, we have $a(t) = (t-1)b(t)$ in $\Z[t]$
    and hence $a \in \im(\sigma-1)$; hence $\ker(N) = \im(\sigma-1)$.
    With the same argument, we can show that $\ker(\sigma-1) = \im(N)$.
\end{proof}

With the property above, we can construct a free resolution of $\Z$ as a $\Z[G]$-module
which is more convenient than the bar resolution.
We call this the 2-\textit{periodic resolution} of $\Z$ as a $\Z[G]$-module.

\begin{corollary}
    The following chain is exact:
    \[
        \cdots \overset{\sigma-1}{\longrightarrow} \Z[G] \overset{N}{\longrightarrow}\Z[G] \overset{\sigma-1}{\longrightarrow} \Z[G] \overset{\epsilon}{\longrightarrow} \Z \to 0
    \]
    where $\epsilon$ is the augmentation map.
\end{corollary}

Suppose now that $A$ is a $G$-module.
Let $B^\bullet \overset{\epsilon}{\longrightarrow} \Z$ be the 2-periodic resolution of $\Z$ as a $\Z[G]$-module.
Then, we have
\begin{align*}
    &H^i(G, A) = H^i(\Hom_{\Z[G]}(B^\bullet, A)),\\
    &\Hom_{\Z[G]}(B^\bullet, A) = \cdots \overset{N}{\longleftarrow} A \overset{\sigma-1}{\longleftarrow} A \overset{N}{\longleftarrow} A \overset{\sigma-1}{\longleftarrow} A.
\end{align*}

Thus, we get
\begin{align*}
    &H^0(G, A) = \ker(\sigma-1) = A^G,\\
    &H^{2i+1}(G, A) = {}_N A/(\sigma-1)A,\\
    &H^{2i+2}(G, A) = A^G/N(A)
\end{align*}
where ${}_N A = \{a \in A: Na = 0\}$ and $i \geq 0$.
We can see that the cohomology groups of a cyclic group are periodic with period 2.
Here are some applications of this computation.

\subsection{Hilbert 90}

Let $L/K$ be a finite Galois extension of fields with Galois group $G = \Gal(L/K)$.
What are the Galois cohomology groups of $G$ with coefficients in $L$ or $L^\times$?
For the $0$-th cohomology group, we can easily see that
$H^0(G, L) = L^G = K$ and $H^0(G, L^\times) = (L^\times)^G = K^\times$.
How about the higher cohomology groups?
\begin{theorem}[Hilbert's Theorem 90]
    $H^1(G, L^\times) = 0$.
    \label{thm_hilbert90}
\end{theorem}

Before proving this theorem, we prove the following lemma.

\begin{lemma}[Linear independence of automorphisms]
    Let $E$ be a field and $\sigma_1,\cdots, \sigma_n$ be distinct automorphisms of $E$.
    Then, for any $c_1,\cdots, c_n \in E$,
    \(
        \sum_{i=1}^n c_i \sigma_i(x) = 0 \text{ for all } x \in E \text{ implies } c_1 = c_2 = \cdots = c_n = 0.
    \)
\end{lemma}
\begin{proof}
    We will prove this by induction on $n$.
    We may assume that none of $c_i$ is zero, otherwise we can remove the corresponding $\sigma_i$ and reduce to a smaller $n$.
    If $n=1$, then we have $c_1 \sigma_1(1) = c_1 \neq 0$.

    Now we consider the case $n>1$ and the statement is true for $n-1$,
    but suppose that $\sum_{i=1}^{n} c_i \sigma_i = 0$ for the sake of contradiction.
    We can set $c_n = -1$ without loss of generality.
    Let $x, y \in E$.
    We have $\sigma_n(x) = \sum_{i=1}^{n-1} c_i \sigma_i(x)$,
    and $\sigma_n(x)\sigma_n(y) = \sigma_n(xy)$.

    The RHS is equal to $\sum_{i=1}^{n-1} c_i \sigma_i(xy) = \sum_{i=1}^{n-1} c_i \sigma_i(x)\sigma_i(y)$
    while the LHS is equal to $\sum_{i=1}^{n-1} c_i \sigma_i(x)\sigma_n(y)$.
    It follows that
    $$0 = \sum_{i=1}^{n-1} c_i \sigma_i(x)\sigma_n(y) - \sum_{i=1}^{n-1} c_i \sigma_i(x)\sigma_i(y) = \sum_{i=1}^{n-1} c_i \sigma_i(x)(\sigma_n(y) - \sigma_i(y)).$$
    By the induction hypothesis, we have $c_i(\sigma_n(y) - \sigma_i(y)) = 0$ for all $i$.
    We also have $c_i \neq 0$ for all $i$, so we must have $\sigma_n(y) - \sigma_i(y) = 0$ for all $i$.
    Recalling that $y$ is arbitrary, we have $\sigma_n = \sigma_i$ for all $i$ which is a contradiction.
\end{proof}

\begin{proof}[Proof of Theorem \ref{thm_hilbert90}]
    Recall that $H^1(G, L^\times)$ is the group of crossed homomorphisms modulo principal crossed homomorphisms.
    Let $\alpha: G \to L^\times$ be a crossed homomorphism.
    For $g, h \in G$, we have $\alpha(gh) = g \alpha(h)\cdot \alpha(g)$.
    We note that each $g \in G$ is an automorphism of $L^\times$, and none of $\alpha(g)$ is zero.
    The lemma implies that there exists $x \in L^\times$ such that $\beta:=\sum_{g \in G} \alpha(g) g(x) \neq 0$.
    For $h \in G$, we have
    \begin{align*}
        h(\beta) &= h\left(\sum_{g \in G} \alpha(g) g(x)\right)
        = \sum_{g \in G} h(\alpha(g)) hg(x)\\
        &= \sum_{g \in G} \alpha(h)^{-1} \alpha(hg) hg(x)
        = \alpha(h)^{-1} \sum_{g \in G} \alpha(hg) hg(x)
        = \alpha(h)^{-1} \beta.
    \end{align*}
    So we have $\alpha(h) = \beta/h(\beta)$ for all $h \in G$.
    Therefore, $\alpha$ is a principal crossed homomorphism, and we conclude that $H^1(G, L^\times) = 0$.
\end{proof}

Let's take a look at the original version of Hilbert 90.
\begin{theorem}[Hilbert's Theorem 90, Original Version]
    Let $L/K$ be a cyclic Galois extension of fields with Galois group $G = \Gal(L/K)$ generated by $\sigma$.
    Let $N: L^\times \to K^\times$ be the norm map given by $N(x) = x\sigma(x)\cdots \sigma^{n-1}(x)$.
    Then, for any $x \in L^\times$ with $N(x) = 1$, there exists $y \in L^\times$ such that $x = y/\sigma(y)$.
\end{theorem}
\begin{proof}
    We want to show that Theorem \ref{thm_hilbert90} implies the original version of Hilbert 90.
    Recall that $H^1(G, L^\times) = \ker(N)/\im(\sigma-1)$ because $G$ is cyclic.
    So if $N(x) = 1$, then $x \in \ker(N)$, and hence $x \in \im(\sigma-1)$ by Theorem \ref{thm_hilbert90}.
    Therefore, there exists $y \in L^\times$ such that $x = y/\sigma(y)$.
\end{proof}

\subsection{Tate Cohomology}
This is the main application of the 2-periodic resolution.
Let $G$ be a cyclic group acting on an abelian group $A$.

\begin{definition}[Tate Cohomology]
    The Tate cohomology groups $\hat{H}^n(G, A)$ are defined as follows:
    \[
        \hat{H}^n(G, A) = \begin{cases}
            A^G/N(A) & n \text{ is even,}\\
            {}_N A/(\sigma-1)A & n \text{ is odd.}
        \end{cases}
    \]
\end{definition}

So now we can consider the negative indices.
In fact, we can do this for groups that are not necessarily cyclic.
We define the negative part with the interaction of cohomology and homology.